\documentclass[10pt,a4paper]{amsart}
\usepackage[margin=19mm]{geometry}
\usepackage{amsmath,amssymb,mathtools}
\usepackage[scr=rsfs,cal=euler]{mathalfa}
\usepackage{xfrac}
\usepackage[unicode,hidelinks,bookmarks=false]{hyperref}
\newcommand{\ie}{i.e.\ }
\newcommand{\cf}{cf.\ }
\newcommand{\resp}{respectively\ }
\newcommand{\Subsection}{\subsection}
\providecommand{\nobreakdash}{\nobreak\hbox{-}}
\setlength{\emergencystretch}{2em}
\multlinegap0pt
\usepackage{algorithm}\usepackage{algpseudocode}
\usepackage{amssymb}
\usepackage{tikz}
\usepackage{bm}
\newtheorem{lemma}{Lemma}
\newcommand{\logb}{\mathrm{\log\text{-}bound}}
\newcommand{\vset}[1]{\mathcal{#1}}
\title{Size Bounds for CQs Under Acyclic Constraints}
\author{Stefan Mengel, Andrei Romashchenko}
\date{Source: arXiv:2608.26775v1 (CC BY 4.0)}
\begin{document}
\maketitle

\noindent\emph{Source and licence.} Size Bounds for CQs Under Acyclic Constraints. Version arXiv:2608.26775v1 (CC BY 4.0), distributed by arXiv under the Creative Commons Attribution 4.0 International licence, http://creativecommons.org/licenses/by/4.0/, as recorded in the arXiv licence metadata for this exact version and shown on the abstract page https://arxiv.org/abs/2608.26775v1. Full licence notice: \texttt{licenses/arxiv-2608.26775-CC-BY-4.0.txt}.\par\smallskip

\noindent\textit{Editorial scope.} Counted unit: the article's lemma lem:tightproject (source0.tex lines 732--738). The article's complete original proof of it is delivered with it (source0.tex lines 741--806, one \texttt{proof} environment of 66 lines). No mathematics is written, repaired or re-derived: the statement and the proof are the article's own lines, in the article's own notation, and no retained line is substituted or re-flowed.  No further source-local context is needed: every cross-reference of every retained line resolves inside the two delivered spans.  Nothing was omitted from any retained span, so the package carries no omission record.  No retained line contains a citation, so the wrapper carries no bibliography and no bibliography entry.  No retained line contains a figure environment, an image-inclusion command or drawing code, so no image is delivered and the package declares no asset dependency.  Editorial formatting only, in addition: an AMS \texttt{amsart} preamble that restates the article's own declarations and macros used by the retained lines, namely the environments \texttt{lemma} on separate counters, together with the standard packages those lines need.  The retained environments print in this document as Lemma 1 in order of appearance, whereas the article prints the same items with the numbering of its own full document.  The complete native source of the article as distributed in the arXiv e-print for this exact version is bundled unchanged as \texttt{sources/208660/source0.tex}.
\begin{lemma}\label{lem:tightproject}
	Let $\vset Y$ be a set of variables, $\vset X\subseteq \vset Y$, $\Pi$ a set of cardinality constraints and $\Sigma$ a set of functional dependencies such that $\vset X$ is a prefix with respect to $\Sigma$. Then
	\begin{align*}
		\logb_{\Gamma^*_{\vset Y}}(\vset X, \Pi, \Sigma) &= \logb_{\Gamma^*_{\vset X}}(\vset X, \Pi^r, \Sigma^r),\\
		\logb_{\Gamma_{\vset Y}}(\vset X, \Pi, \Sigma) &= \logb_{\Gamma_{\vset X}}(\vset X, \Pi^r, \Sigma^r).
	\end{align*}
\end{lemma}

\begin{proof}
\emph{Claim 1:}	 $\logb_{\Gamma_{\vset Y}}(\vset X, \Pi, \Sigma) \ge \logb_{\Gamma_{\vset X}}(\vset X, \Pi^r, \Sigma^r)$. 

Consider any polymatroid $h$ on $\vset X$ satisfying $\Pi^r$ and $\Sigma^r$. We will construct a polymatroid $h'$ that satisfies $\Pi$ and $\Sigma$ with $h(\vset X)= h'(\vset X)$. From this, the claim follows directly. To construct $h'$, we extend $h$ to a polymatroid on $\vset Y$ by setting for every $\vset A\subseteq Y$ the value $h'(\vset A) := h(\vset A\cap \vset X)$. Clearly, $h'$ is a polymatroid satisfying $\Pi$. We show that it also satisfies the constraints in $\Sigma$. To this end, consider an FD $\vset W\rightarrow \vset U$ from $\Sigma$. We claim that
\[
h'(\vset U\cup \vset W) = h((\vset U\cup \vset W)\cap \vset X) = h((\vset U\cap \vset X)\cup (\vset W\cap \vset X)) = h(\vset W\cap \vset X) = h'(\vset W).
\] 
In this chain of equalities only the third one is non-trivial. In the case $\vset U \cap \vset X = \emptyset$ it is straightforward; 
and in the other case it
holds because $\vset W \cap \vset X \rightarrow \vset U\cap \vset X$ is an FD in $\Sigma^r$ and thus satisfied by $h$. 
Since $h'(\vset X) = h(\vset X)$, it follows that $\logb_{\Gamma_{\vset Y}}(\vset X, \Pi, \Sigma) \ge \logb_{\Gamma_{\vset X}}(\vset X, \Pi^r, \Sigma^r)$, as claimed.


\smallskip
	 
\emph{Claim 2:} $\logb_{\Gamma^*_{\vset Y}}(\vset X, \Pi, \Sigma) \ge \logb_{\Gamma^*_{\vset X}}(\vset X, \Pi^r, \Sigma^r)$. 

To prove this claim, consider any entropic function $h$ on $\vset X$ that satisfies all constraints in $\Pi^r$ and $\Sigma^r$. We construct a new entropic function $h'$ on $\vset Y$ by adding random variables for $\vset Y\setminus \vset X$ and letting their distribution be such that they only ever take a fixed value. Then for any subset $\vset A\subseteq \vset Y\setminus \vset X$ we have $h'(\vset A) = H(\vset Y_{\vset A}) = 0$ by definition of the entropy $H$. Since $h'$ is entropic and thus a polymatroid, it follows that for every $\vset A\subseteq \vset Y$ that $h'(\vset A\cap \vset X) \le h'(\vset A) \le h'(\vset A\cap \vset X) + h'(\vset A\setminus \vset X) = h'(\vset A\cap \vset X)$, so $h'(\vset A) = h'(\vset A \cap \vset X) = h(\vset A\cap \vset X)$. Now the claim follows exactly as for the polymatroid case.
	
\smallskip
	 
\emph{Claim 3:}
$\logb_{\Gamma_{\vset Y}}(\vset X, \Pi, \Sigma) \le \logb_{\Gamma_{\vset X}}(\vset X, \Pi^r, \Sigma^r)$.

To prove this claim, we consider a polymatroid $h\in \Gamma_{\vset Y}$ satisfying the constraints in $\Pi$ and $\Sigma$. We define $h': \vset Y\rightarrow \mathbb{R}$ by setting  $h'(\vset A) = h(\vset A\cap \vset X)$ for every $\vset A\subseteq \vset Y$.
%
%
We claim that $h'$ defined this way is a polymatroid.)
Clearly, $h'(\emptyset) = 0$ and $h'$ is monotone, so it only remains to show that $h'$ is submodular. So let $\vset A, \vset B\subseteq \vset Y$. Then 
	\begin{align*}
		h'(\vset A \cap \vset B) + h'(\vset A \cup \vset B) & = h((\vset A \cap \vset B)\cap \vset X) + h((\vset A \cup \vset B)\cap \vset X)\\
		&= h((\vset A \cap \vset X)\cap (\vset B\cap \vset X)) + h((\vset A \cap \vset X) \cup (\vset B\cap \vset X))\\
		&\le h(\vset A \cap \vset X) + h(\vset B\cap \vset X)%
		%
		= h'(\vset A) + h'(\vset B),
	\end{align*}
	which shows that $h'$ is indeed a polymatroid.
	
	We now claim that $h'$ satisfies the FDs in $\Sigma'$. So consider an FD $\vset W\cap \vset X \rightarrow \vset U\cap \vset X$ from $\Sigma'$. Assume first that $\vset W\subseteq \vset X$. Then 
	\begin{align*}
		h'(\vset W\vset U) - h'(\vset W) & = h((\vset W \cup \vset U)\cap \vset X) - h(\vset W)%
		%
		\le h(\vset W\cup \vset U) - h(\vset W)%
		%
		= 0,
	\end{align*}
	where the inequality holds because $h$ is monotone and the final equality it true because $h$ satisfies $\Sigma$ and thus $\vset W \rightarrow \vset U$. Now consider the case that $\vset W \nsubseteq \vset X$. Then, because $\vset X$ is a prefix, we have that $\vset U \cap \vset X = \emptyset$. It follows that 
	\begin{align*}
	h'(\vset W\vset U) - h'(\vset W) & = h((\vset W \cup \vset U)\cap \vset X) - h(\vset W\cap \vset X)%
	%
	= h(\vset W \cap \vset X) - h(\vset W\cap \vset X)%
	%
	= 0.
	\end{align*}
	So in any case $h'$ satisfies $\vset W\cap \vset X \rightarrow \vset U \cap \vset X$ and thus $\Sigma'$.
	
	We now define a set function $h''$ on $\vset X$ by restricting $h'$ to the ground set $\vset X$. Then $h''$ is a polymatroid that satisfies the constraints in $\Sigma'$. Moreover, we have $h''(\vset X) = h'(\vset X) = h(\vset X)$. So for every polymatroid $h$ on $\vset X$ respecting $\Sigma$, we can construct a polymatroid $h''$ on $\vset X$ respecting $\Sigma'$ such that $h''(\vset X) = h(\vset X)$, and it follows that $\logb_{\Gamma_{\vset Y}}(\vset X, \Pi, \Sigma) \le \logb_{\Gamma_{\vset X}}(\vset X, \Pi^r, \Sigma^r)$.

\smallskip
	 
\emph{Claim 4:}	
	 $\logb_{\Gamma^*_{\vset Y}}(\vset X, \Pi, \Sigma) \le \logb_{\Gamma^*_{\vset X}}(\vset X, \Pi^r, \Sigma^r)$.
	 
	 To prove this claim, we  let $h$ be an entropic function on $\vset Y$. Assume for simplicity that $\vset Y = \{1, \ldots, n\}$. Then there are random variables $\{X_1 , \ldots, X_n\}$ such that for every $\mathcal A\subseteq \vset Y$ we have $h(A) = H(X_{\vset A})$ where $X_{\vset A} = \{X_i\mid i\in \vset A\}$. Then define a set of random variables $\{X_1', \ldots, X_n'\}$ having the same distribution for the $\{X'_i\mid i\in \vset X\}$ as for $\{X_i\mid i\in \vset X\}$ and fixing all other variables to a constant value. Let $h'$ be the entropic function induced by the set $\{X_1', \ldots, X_n' \}$. Then we have for every $\vset B \subseteq \vset Y \setminus\vset X$ that $h'(\vset B) = 0$, since the corresponding random variables are constant. Since $h'$ is entropic and thus a polymatroid, it follows that for every $\vset A \subseteq \vset Y$ we have $h'(\vset A \cap \vset X) \le h'(\vset A) \le  h'(\vset A \cap \vset X) + h'(\vset A \setminus \vset X) \le h'(\vset A \cap \vset X)$ and thus $h'(\vset A) = h'(\vset A \cap \vset X)$. However, by definition of $h'$, we have $h'(\vset A\cap \vset X) = h(\vset A\cap \vset X)$ since both values are the entropy of the same random variable $X_{\vset A\cap \vset X}$. But then $h'(\vset A) = h(\vset A\cap \vset X)$. Thus, $h'$ is exactly as in 
	 Claim~3 (the case of polymatroids) but here has the additional property that it is entropic. The rest of the argument follows as the polymatroid case.
\end{proof}

\end{document}
